\scanpage{10}

\[\resizebox{\ifdim\width>\linewidth\linewidth\else\width\fi}{!}{$\displaystyle
 \boxed{\mathcal F_1\subset\Omega_{\deleted{\unsure{1}}}}
 \quad\deleted{\unsure{\varnothing}}\quad \text{i.e.}\quad
 \mathcal F_1\in{\unsure{\mathscr P|\wp|\mathfrak P}}^{\ast}(\Omega)
$}\]
\[
 \widetilde{\mathcal F}_1=(\widetilde\Omega\mid\mathcal F_1)\cap T_1
 \quad\subset\widetilde\Omega=\widetilde\Omega_1
\]
\[
 \boxed{f_2:\widetilde{\mathcal F}_1\longrightarrow\mathscr P^\ast(\Omega)}
 \qquad\widetilde\Omega\times\Omega
\]

$\deleted{\unsure{\widetilde{\mathcal F}_1}}\
(\widetilde\omega_1,\omega_2)\in\widetilde\Omega\times\Omega$ comp. avec $(\mathcal F_1,\mathcal F_2)$ si
\[\resizebox{\ifdim\width>\linewidth\linewidth\else\width\fi}{!}{$\displaystyle
\begin{aligned}
 &\text{a)}\quad\widetilde\omega_1\in\widetilde{\mathcal F}_1
 \quad\text{i.e.}\quad\omega_1\in\mathcal F_1\text{ et }\widetilde\omega_1\in T_1\\
 &\deleted{\text{b)}\quad\widetilde\omega_1\in T_1}\\
 &\text{b)}\quad\omega_2\in f_2(\widetilde\omega_1)
\end{aligned}
$}\]
\[\resizebox{\ifdim\width>\linewidth\linewidth\else\width\fi}{!}{$\displaystyle
\begin{aligned}
 \widetilde{\mathcal F}_2\subset\widetilde\Omega_2
 =\widetilde\Omega\times\widetilde\Omega\ \unsure{\ni|=}\ 
 \bigl\{(\widetilde\omega_1,\omega_2)\bigm|&
 \omega_1\in\mathcal F_1,\,
 (\widetilde\omega_1\times\widetilde\omega_2)\in T_2,\\
 &\widetilde\omega_2\in f_2(\widetilde\omega_1)\bigr\}
\end{aligned}
$}\]
\[
 f_3:\widetilde{\mathcal F}_2\longrightarrow\mathscr P(\Omega)
\]
\[\resizebox{\ifdim\width>\linewidth\linewidth\else\width\fi}{!}{$\displaystyle
\begin{aligned}
 &\cdots\\
 &f_\nu:\underset{\substack{\cap\\\widetilde\Omega^{\nu-1}}}{\widetilde{\mathcal F}_{\nu-1}}
 \longrightarrow\mathscr P(\Omega)\qquad\text{\unsure{d'où|déf|donc}}\\
 &\widetilde{\mathcal F}_\nu\subset\widetilde\Omega^\nu\quad
 \text{(en fait, }\widetilde{\mathcal F}_\nu\subset T_\nu)
\end{aligned}
$}\]

$(\mathcal F_1,\ldots,\mathcal F_\nu)$ est appelé \unsure{une|un} \deleted{\unsure{multiplicité|multiplet} de} \inserted{\text{crible de \uuline{longueur $\nu$}}}

\uline{NB} \unsure{certains|Certains} $\widetilde{\mathcal F}_i$ peuvent \unsure{être|etre} $\varnothing$

\[\resizebox{\ifdim\width>\linewidth\linewidth\else\width\fi}{!}{$\displaystyle
\begin{aligned}
 \mathcal F_1&=f_1\subset\Omega
 &\widetilde\Omega\times\widetilde\Omega\ \cdots\ \times\widetilde\Omega\\
 \mathcal F_2&\subset\widetilde\Omega\times\Omega&\mathcal F\\
 \mathcal F_3&\subset\widetilde\Omega\times\widetilde\Omega\times\Omega\\
 &\cdots\\
 \mathcal F_\nu&\subset
 \underbrace{\widetilde\Omega\times\widetilde\Omega\ \cdots\ \widetilde\Omega}_{\nu-1}\times\Omega
\end{aligned}
$}\]
% archive sheet number: 16
